\documentclass[11pt,a4paper]{article}
\usepackage[utf8]{inputenc}
\usepackage[T1]{fontenc}
\usepackage[spanish]{babel}
\usepackage{amsmath,amssymb,mathtools}
\usepackage{geometry}
\usepackage{enumitem}
\usepackage{lmodern}
\usepackage{microtype}
\geometry{margin=2.4cm}
\setlength{\parindent}{0pt}
\setlength{\parskip}{0.5em}
\newcommand{\R}{\mathbb{R}}
\newcommand{\Span}{\operatorname{L}}
\newcommand{\Ima}{\operatorname{Im}}
\newcommand{\sen}{\operatorname{sen}}
\begin{document}

\begin{center}
{\LARGE\bfseries Álgebra Lineal}\\[0.4em]
{\Large Tema 4: El Espacio Afín --- Soluciones}
\end{center}

\section{Ejercicio 1}
\subsection{Solución}

\[
\overrightarrow{OP}
=
\overrightarrow{OO'}+\overrightarrow{O'P}
\]

\[
x_1u_1+x_2u_2
=
2u_1+3u_2+y_1(u_1-u_2)+y_2(5u_1+4u_2)
\]

Por tanto:

\[
\boxed{
x_1=2+y_1+5y_2,
\qquad
x_2=3-y_1+4y_2
}
\]

\[
\boxed{
\begin{pmatrix}
1\\x_1\\x_2
\end{pmatrix}
=
\begin{pmatrix}
1&0&0\\
2&1&5\\
3&-1&4
\end{pmatrix}
\begin{pmatrix}
1\\y_1\\y_2
\end{pmatrix}
}
\]

\section{Ejercicio 2}
\subsection{Solución}

Sea $P_1$ una solución particular:

\[
AP_1=C.
\]

Si $P$ es otra solución:

\[
AP=C.
\]

Restando:

\[
A(P-P_1)=0.
\]

Por tanto:

\[
P-P_1\in\Ker A
\]

y:

\[
\boxed{
\{X:AX=C\}=P_1+\Ker A.
}
\]

\section{Ejercicio 3}
\subsection{Solución}

Igualando las ecuaciones de $F_1$ y $F_2$:

\[
a+3\lambda+2\mu=2+\alpha+2\beta
\]

\[
1-\lambda-\mu=1
\]

\[
4+\lambda=1+\alpha+\beta
\]

\[
6+5\lambda+2\mu=3\alpha.
\]

Resolviendo:

\[
\boxed{a=6}
\]

Un punto común es:

\[
P=(6,1,4,6).
\]

Los espacios de dirección son:

\[
V_1=
\Span\{(3,-1,1,5),(2,-1,0,2)\}
\]

\[
V_2=
\Span\{(1,0,1,3),(2,0,1,0)\}.
\]

Se obtiene:

\[
V_1\cap V_2=
\Span\{(1,0,1,3)\}.
\]

Por tanto:

\[
\boxed{
F_1\cap F_2
=
(6,1,4,6)+\Span\{(1,0,1,3)\}
}
\]

\[
\boxed{
F_1+F_2
=
(6,1,4,6)
+
\Span\{(2,-1,0,2),(2,0,1,0),(3,-1,1,5)\}.
}
\]

\section{Ejercicio 4}
\subsection{Solución}

\[
\overrightarrow{AB}=(-1,1,2)
\]

\[
\overrightarrow{AC}=(-2,-1,1)
\]

\[
\overrightarrow{AD}=(1,2,0).
\]

Como:

\[
\overrightarrow{OP}
=
\overrightarrow{OA}+\overrightarrow{AP},
\]

\[
\begin{pmatrix}
x\\y\\z
\end{pmatrix}
=
\begin{pmatrix}
3\\1\\-2
\end{pmatrix}
+
\begin{pmatrix}
-1&-2&1\\
1&-1&2\\
2&1&0
\end{pmatrix}
\begin{pmatrix}
x'\\y'\\z'
\end{pmatrix}.
\]

Por tanto:

\[
\boxed{
x=3-x'-2y'+z',
\quad
y=1+x'-y'+2z',
\quad
z=-2+2x'+y'.
}
\]

\section{Ejercicio 5}
\subsection{Solución}

Matricialmente:

\[
\boxed{
\begin{pmatrix}
x_1\\x_2\\x_3
\end{pmatrix}
=
\begin{pmatrix}
1\\0\\2
\end{pmatrix}
+
\begin{pmatrix}
0&1&0\\
\frac13&0&-\frac23\\
\frac23&0&\frac13
\end{pmatrix}
\begin{pmatrix}
y_1\\y_2\\y_3
\end{pmatrix}
}
\]

Sustituyendo en:

\[
x_1+x_2+x_3=3
\]

se obtiene:

\[
\boxed{
3y_1+3y_2-y_3=0.
}
\]

En $R$:

\[
\boxed{
\pi=
(0,0,3)
+
\Span\{(1,0,-1),(0,1,-1)\}.
}
\]

En $R'$:

\[
\boxed{
\pi=
(0,0,0)
+
\Span\{(1,0,3),(0,1,3)\}.
}
\]

\section{Ejercicio 6}
\subsection{Solución}

\[
F_1=P_1+V_1
\]

con:

\[
P_1=(0,4,0,a),
\]

\[
V_1=
\Span\{(1,-1,0,0),(0,0,1,-1)\}.
\]

Además:

\[
F_2=P_2+V_2
\]

con:

\[
P_2=(3,2,0,-1),
\]

\[
V_2=\Span\{(1,-2,2,1)\}.
\]

La dimensión es mínima cuando:

\[
\overrightarrow{P_1P_2}
=
(3,-2,0,-1-a)
\]

pertenece a $V_1+V_2$.

La condición da:

\[
\boxed{a=2}
\]

y:

\[
\boxed{
F_1+F_2
=
(0,4,0,2)
+
\Span\{(1,-1,0,0),(0,0,1,-1),(1,-2,2,1)\}.
}
\]

\section{Ejercicio 7}
\subsection{Solución}

El polinomio:

\[
x
\]

pertenece a $V$.

El espacio de dirección viene dado por:

\[
p(1)=0,
\qquad
p(2)=0.
\]

Por tanto:

\[
\boxed{
V=x+\Span\{(x-1)(x-2)\}.
}
\]

Un sistema de referencia afín es:

\[
\boxed{
R_V=\{x;(x-1)(x-2)\}.
}
\]

Además:

\[
2x^2-5x+4
=
x+2(x-1)(x-2).
\]

Así:

\[
\boxed{
[2x^2-5x+4]_{R_V}=2.
}
\]

El vector director en la base canónica es:

\[
(2,-3,1).
\]

El hiperplano perpendicular que pasa por $x=(0,1,0)$ verifica:

\[
\boxed{
2a-3b+c=-3.
}
\]

\section{Ejercicio 8}
\subsection{Solución}

\[
M(\widetilde f)
=
\begin{pmatrix}
1&1\\
4&-1\\
0&1
\end{pmatrix}.
\]

Sea:

\[
f(X)=C+M(\widetilde f)X.
\]

Como:

\[
f(1,2)=(1,2,3),
\]

\[
C=
\begin{pmatrix}
1\\2\\3
\end{pmatrix}
-
\begin{pmatrix}
3\\2\\2
\end{pmatrix}
=
\begin{pmatrix}
-2\\0\\1
\end{pmatrix}.
\]

Por tanto:

\[
\boxed{
\begin{pmatrix}
y_1\\y_2\\y_3
\end{pmatrix}
=
\begin{pmatrix}
-2\\0\\1
\end{pmatrix}
+
\begin{pmatrix}
1&1\\
4&-1\\
0&1
\end{pmatrix}
\begin{pmatrix}
x_1\\x_2
\end{pmatrix}
}
\]

y:

\[
\boxed{
\Ima f
=
(-2,0,1)
+
\Span\{(1,4,0),(1,-1,1)\}.
}
\]

\section{Ejercicio 9}
\subsection{Solución}

El cambio es:

\[
\begin{pmatrix}
x_1\\x_2
\end{pmatrix}
=
\begin{pmatrix}
3\\3
\end{pmatrix}
+
\begin{pmatrix}
\frac1{\sqrt2}&\frac2{\sqrt2}\\
\frac1{\sqrt2}&-\frac2{\sqrt2}
\end{pmatrix}
\begin{pmatrix}
y_1\\y_2
\end{pmatrix}.
\]

Sustituyendo en:

\[
(x_1-1)^2+(x_2-1)^2=9
\]

se obtiene:

\[
\boxed{
(y_1+2\sqrt2)^2+4y_2^2=9.
}
\]

Por tanto:

\[
\boxed{
C\text{ es una elipse en }R'.
}
\]

\section{Ejercicio 10}
\subsection{Solución}

\[
\begin{pmatrix}
x_1\\x_2\\x_3
\end{pmatrix}
=
\begin{pmatrix}
1\\1\\0
\end{pmatrix}
+
\begin{pmatrix}
-1&0&1\\
0&-1&1\\
1&1&1
\end{pmatrix}
\begin{pmatrix}
x_1'\\x_2'\\x_3'
\end{pmatrix}.
\]

Así:

\[
x_1=1-x_1'+x_3'
\]

\[
x_2=1-x_2'+x_3'
\]

\[
x_3=x_1'+x_2'+x_3'.
\]

Sustituyendo:

\[
x_1+x_2+x_3=2
\]

se obtiene:

\[
\boxed{x_3'=0.}
\]

\section{Ejercicio 11}
\subsection{Solución}

Para $A_1$ y $A_2$ tomamos:

\[
V_1=\Span\{(1,1,1)\},
\]

\[
V_2=\Span\{(2,1,2)\}.
\]

El conjunto:

\[
\{(1,1,1),(2,1,2),(1,0,0)\}
\]

tiene rango $3$.

Por tanto:

\[
\boxed{\dim(A_1+A_2)=3}
\]

y:

\[
\boxed{
\dim(A_1+A_2)=\dim A_1+\dim A_2+1.
}
\]

Para:

\[
A_3=\{(x,y,z):2x-y-z=0\},
\]

\[
\dim A_3=2.
\]

Se obtiene:

\[
\boxed{\dim(A_1+A_3)=3}
\]

y:

\[
\boxed{
\dim(A_1+A_3)=\dim A_1+\dim A_3.
}
\]

\section{Ejercicio 12}
\subsection{Solución}

Toda aplicación de la forma $C+AX$ con $A$ lineal es afín.

\[
\boxed{\text{a) Afín}}
\]

\[
\boxed{\text{b) Afín}}
\]

\[
\boxed{\text{c) Afín}}
\]

\[
\boxed{\text{d) Afín}}
\]

En e):

\[
f(x,y)=(x^2+1,y)
\]

contiene un término cuadrático.

\[
\boxed{\text{e) No es afín}}
\]

\section{Ejercicio 13}
\subsection{Solución}

La matriz del giro de ángulo $\pi/2$ es:

\[
A=
\begin{pmatrix}
0&-1\\
1&0
\end{pmatrix}.
\]

El centro $(1,1)$ es fijo, por lo que:

\[
b=
\begin{pmatrix}
1\\1
\end{pmatrix}
-
A
\begin{pmatrix}
1\\1
\end{pmatrix}
=
\begin{pmatrix}
2\\0
\end{pmatrix}.
\]

Así:

\[
\boxed{
\begin{pmatrix}
x'\\y'
\end{pmatrix}
=
\begin{pmatrix}
2\\0
\end{pmatrix}
+
\begin{pmatrix}
0&-1\\
1&0
\end{pmatrix}
\begin{pmatrix}
x\\y
\end{pmatrix}
}
\]

\[
\boxed{
x'=2-y,
\qquad
y'=x.
}
\]

\section{Ejercicio 14}
\subsection{Solución}

Un vector director de:

\[
x+y=1
\]

es:

\[
d=(1,-1).
\]

La proyección lineal sobre $\Span\{d\}$ es:

\[
A=
\frac12
\begin{pmatrix}
1&-1\\
-1&1
\end{pmatrix}.
\]

Tomando el punto fijo $(1,0)$:

\[
b=
\begin{pmatrix}
1\\0
\end{pmatrix}
-
A
\begin{pmatrix}
1\\0
\end{pmatrix}
=
\begin{pmatrix}
\frac12\\[1mm]
\frac12
\end{pmatrix}.
\]

Por tanto:

\[
\boxed{
x'=\frac{1+x-y}{2},
\qquad
y'=\frac{1-x+y}{2}.
}
\]

\section{Ejercicio 15}
\subsection{Solución}

La matriz de la simetría lineal respecto de la dirección $(1,-1)$ es:

\[
A=
\begin{pmatrix}
0&-1\\
-1&0
\end{pmatrix}.
\]

Tomando $(1,0)$ como punto fijo:

\[
b=
\begin{pmatrix}
1\\0
\end{pmatrix}
-
A
\begin{pmatrix}
1\\0
\end{pmatrix}
=
\begin{pmatrix}
1\\1
\end{pmatrix}.
\]

Así:

\[
\boxed{
x'=1-y,
\qquad
y'=1-x.
}
\]

\section{Ejercicio 16}
\subsection{Solución}

Un punto de $F$ es:

\[
P_0=(1,0,0).
\]

El espacio de dirección es:

\[
S=
\{(x_1,x_2,x_3):x_1+x_2-x_3=0\}.
\]

Como:

\[
x_3=x_1+x_2,
\]

\[
S=
\Span\{(1,0,1),(0,1,1)\}.
\]

Por tanto:

\[
\boxed{
F=
(1,0,0)
+
\Span\{(1,0,1),(0,1,1)\}
}
\]

y:

\[
\boxed{\dim F=2.}
\]

En el nuevo sistema:

\[
\begin{pmatrix}
x_1\\x_2\\x_3
\end{pmatrix}
=
\begin{pmatrix}
-2\\1\\3
\end{pmatrix}
+
\begin{pmatrix}
1&0&1\\
0&1&1\\
-1&-1&1
\end{pmatrix}
\begin{pmatrix}
y_1\\y_2\\y_3
\end{pmatrix}.
\]

Entonces:

\[
x_1=-2+y_1+y_3
\]

\[
x_2=1+y_2+y_3
\]

\[
x_3=3-y_1-y_2+y_3.
\]

Sustituyendo en:

\[
x_1+x_2-x_3=1
\]

se obtiene:

\[
\boxed{
2y_1+2y_2+y_3=5.
}
\]

\end{document}
